\documentclass[11pt,a4paper]{article}

\usepackage[a4paper, top=2.2cm, bottom=2.2cm, left=2.4cm, right=2.4cm]{geometry}
\usepackage[T1]{fontenc}
\usepackage{lmodern}
\usepackage{microtype}
\usepackage{xcolor}
\usepackage{array}
\usepackage{enumitem}
\usepackage{titlesec}
\usepackage[hidelinks]{hyperref}

\definecolor{forest}{HTML}{2D7355}
\definecolor{deepforest}{HTML}{1B3F2F}
\definecolor{ink}{HTML}{14181A}
\definecolor{rule}{HTML}{1B3F2F}

\titleformat{\section}
  {\color{forest}\sffamily\large\bfseries}{\thesection}{0.6em}{}
\titlespacing*{\section}{0pt}{14pt}{5pt}

\setlist{topsep=3pt, itemsep=2pt, parsep=0pt}
\setlength{\parskip}{0.4em}
\setlength{\parindent}{0pt}
\color{ink}

\begin{document}
\thispagestyle{empty}

\noindent{\sffamily\color{forest}\bfseries\small PROJECT MEMO \quad$\cdot$\quad CHESS-NN-PLAYGROUND}

\vspace{0.3cm}
\noindent{\sffamily\bfseries\LARGE\color{deepforest} Reframing the project into a research paper}

\vspace{0.15cm}
\noindent{\itshape\color{forest}\large A response to supervisor feedback, with two questions for guidance.}

\vspace{0.25cm}
\noindent{\color{rule}\rule{\linewidth}{1.2pt}}

\vspace{0.2cm}
\noindent
\begin{tabular}{@{}>{\sffamily\bfseries\footnotesize}l@{\hspace{16pt}}>{\footnotesize}l@{}}
AUTHOR & Lennart Axel Conrad (Student ID: 2025080264) \\
AFFILIATION & Zijing College, Tsinghua University \\
SUPERVISOR & Prof.\ Jungong Han, Department of Automation \\
SUBJECT & Reframing the chess-architecture investigation into a contribution-driven paper \\
\end{tabular}

\section{Where the project stands}

The current article evaluates a broad range of neural architectures for chess-position
classification, including meaningful modifications of existing designs. I agree with your
assessment: at present it reads as a strong technical investigation rather than a research
paper with clearly defined, validated contributions. This memo proposes a sharper framing
and asks for your guidance on two decisions I cannot make alone.

\section{Proposed contributions}

\textbf{Contribution 1 --- A hard-negative chess-position benchmark.} \textit{This is the
main contribution.} Ordinary puzzle classification is too easy; the real difficulty comes
from \emph{near-puzzles} --- positions that look tactical but are not verified
unique-solution puzzles. I propose a puzzle-vs-near-puzzle benchmark, evaluated with a
matched-recall near-puzzle false-positive metric, that exposes a representation-failure mode
which generic architectures and chess-specific baselines all struggle with.

\textbf{Contribution 2 --- A new architecture for puzzle classification, and a mechanism
study of its key design choice.} My broad architecture scout served as \emph{hypothesis
generation}, not as final evidence. Its clearest signal is a chess-specific architecture
built around one primitive: a dual-stream design that separates material/exchange tactics
from king-safety pressure, combined by a learned router. This has an intuitive justification
(it mirrors how chess evaluation is classically decomposed) and a testable prediction (it
should specifically reduce near-puzzle false positives). The paper would present the
architecture and then validate the decomposition itself through per-component ablations,
multiple seeds, and a scaling study.

\section{Honest status of the current evidence}

Most experiments so far are single-seed, at a single model size, and on a subset of the
data --- this is screening evidence, not scientific proof. In particular, my most promising
candidate architecture looks strong at small scale, but I do not yet know whether the
advantage survives scaling. I will treat that as a question to answer, not a claim to make.

\section{Two questions I would value your guidance on}

\textbf{1. Direction of the project.} Everything is currently framed as classification.
I see four possible directions, ordered roughly from most contained to most ambitious, and
I am genuinely unsure which makes the strongest paper:
\begin{itemize}
  \item \textbf{(a) Stay with classification} --- keep the hard-negative benchmark as the
  target; the contribution is a rigorous benchmark + architecture + ablation study. Lowest
  risk, and achievable on my current hardware.
  \item \textbf{(b) Move to tactical move selection} --- keep the same near-puzzle idea,
  but ask a more chess-native question: can the model identify or rank the tactical move,
  not just classify the position as puzzle-like? This is between classification and a full
  engine. It would test whether the architecture learns actionable tactical structure
  rather than a binary surface signal, while still being much cheaper than full engine
  training.
  \item \textbf{(c) Move toward a chess engine} --- replace the classification head with
  policy and win--draw--loss heads, as real engines use, and evaluate the architecture in
  actual play or fixed-depth search against a baseline. This is the most ambitious
  chess-specific direction and the most practically exciting, but it is also the most
  data- and compute-hungry.
  \item \textbf{(d) Abstract the design beyond chess} --- treat the exchange/king split
  as one example of a broader principle: structured task decomposition with a learned
  router for hard-negative reasoning problems. It would be scientifically valuable to test
  whether this transfers to another structured domain with confusable positives and
  negatives, such as theorem-proving states, program-analysis tasks, graph reasoning, or
  board-game problems outside chess. This has the highest scientific ceiling because the
  contribution would no longer be limited to chess, but it broadens the paper and would
  require a second testbed.
\end{itemize}
My instinct is that the paper's spine should be (a), with a careful mechanism study of the
best architecture. Direction (b) is an attractive middle ground if classification alone
feels too narrow; (c) is the long-term chess-engine goal; and (d) is the highest-level
scientific framing if we want the contribution to generalise beyond chess. I would value
your opinion on the right scope, especially given the compute question below.

\textbf{2. Compute resources.} This is my main practical constraint. My current setup is a
single GPU on a remote machine --- enough for small-scale screening, but not for the
rigorous validation the paper needs. Multiple seeds, multiple model sizes, a scaling study,
and a from-scratch baseline comparison would take weeks of wall-clock time on my current
hardware. If there is any way to access additional compute through the department or the
lab, it would be the difference between a thorough validation and a partial one, and I would
be very grateful for help arranging that.

\section{Planned validation steps}

\begin{enumerate}
  \item Lock one canonical dataset split and reconcile inconsistent reports.
  \item Re-run the finalist architectures and baselines with multiple seeds.
  \item Run controlled component ablations of the exchange/king decomposition.
  \item Run a scaling study --- on the finalist architectures only.
  \item Compare against a BT4-style architecture trained from scratch under identical data,
  optimiser, schedule, and compute budget.
  \item Report an efficiency analysis (parameters, FLOPs, inference latency) as
  accuracy/robustness-vs-cost Pareto fronts.
\end{enumerate}

\vspace{0.3cm}
\noindent{\color{rule}\rule{\linewidth}{0.6pt}}

\vspace{0.15cm}
\noindent\textit{\color{forest} If this direction sounds reasonable, I would be glad to meet
briefly to discuss priorities --- especially the compute question, which gates how much of
the validation plan is feasible before the deadline.}

\end{document}
